\subsection{1. Đặt vấn đề}

\indent\indent Trong thời đại số, các giá trị của di sản văn hóa phi vật thể (ICH) ngày càng được tôn vinh và lan tỏa rộng rãi nhờ tính tiếp cận thuận lợi và nhanh chóng thông qua các nền tảng số, mạng xã hội và các phương tiện truyền thông khác. Việc giữ gìn, bảo tồn và phát huy các giá trị này không chỉ giúp bảo vệ bản sắc văn hóa dân tộc mà còn mở ra cơ hội quảng bá, 
nghiên cứu và ứng dụng vào đời sống xã hội, du lịch, giáo dục và các ngành công nghiệp sáng tạo. Văn hóa phi vật thể bao gồm nhiều loại hình khác nhau như lễ hội, nghệ thuật biểu diễn, nghề thủ công truyền thống, phong tục tập quán và tri thức dân gian, mỗi loại hình lại gắn với các yếu tố đặc trưng về địa điểm, thời gian, người thực hiện và các hiện vật liên quan.\vspace{0.3cm}

Đặc điểm của văn hóa phi vật thể là tính đa phương thức và phức tạp. Mỗi thực thể văn hóa thường tồn tại đồng thời dưới nhiều dạng dữ liệu: hình ảnh lễ hội, video trình diễn, âm thanh nhạc cụ, bài hát dân gian, văn bản mô tả lịch sử và tri thức, cùng với các mối quan hệ cấu trúc giữa các thực thể như địa điểm, nghệ nhân, loại hình nghệ thuật, hay hiện vật đi kèm. 
Bên cạnh đó, các yếu tố vật chất, trang phục, nghi thức hay hiện vật thường có sự giao thoa và phụ thuộc lẫn nhau, làm tăng độ khó trong việc phân loại nếu chỉ dựa vào một phương thức dữ liệu đơn lẻ. Những hạn chế này khiến các phương pháp truyền thống như MLP, CNN đơn phương thức bỏ sót các mối liên hệ ẩn trong dữ liệu.
\vspace{0.3cm}

Trước những thách thức trên, tôi đề xuất phương án sử dụng đồ thị đa phương thức (Multimodal Graph -- MG) kết hợp GNN và KAN. Với Multimodal Graph Neural Network (MGNN) sẽ khai thác đồng thời đặc trưng của từng loại dữ liệu và các mối quan hệ giữa các thực thể văn hóa, cho phép mô hình hiểu sâu các ngữ cảnh nội tại cũng như liên phương thức (hình ảnh, văn bản, âm thanh). Tiếp đó, mạng KAN được sử dụng để mô hình hóa các quan hệ phi tuyến phức tạp, tìm ra hàm xấp xỉ tối ưu trong không gian đặc trưng đa chiều, tăng khả năng khái quát hóa của mô hình. Sự kết hợp này kỳ vọng nâng cao hiệu quả phân loại, vượt trội so với các mô hình truyền thống chỉ dựa vào một modal, đồng thời đóng góp vào việc bảo tồn, số hóa và khai thác tri thức văn hóa phi vật thể.